\subsection{The Null-Sentinel Object (\texttt{None})}

Python uses \texttt{None} to represent absence, missing result, or ``no useful value here''. It is not the same thing as zero, an empty string, an empty list, or \texttt{False}.

A first inspection:

\begin{verbatim}
value = None

print(f"value: {value}")
print(f"type:  {type(value)}")
print(f"id:    {id(value)}")

print()
print("selected None methods:")
print(f"{'__bool__':<15} {'__bool__' in dir(value)}")
print(f"{'__repr__':<15} {'__repr__' in dir(value)}")
print(f"{'__eq__':<15} {'__eq__' in dir(value)}")
print(f"{'__class__':<15} {'__class__' in dir(value)}")

print()
print(f"bool(None): {bool(None)}")
\end{verbatim}

Possible output:

\begin{verbatim}
value: None
type:  <class 'NoneType'>
id:    140734431489696

selected None methods:
__bool__        True
__repr__        True
__eq__          True
__class__       True

bool(None): False
\end{verbatim}

The type of \texttt{None} is \texttt{NoneType}. The value \texttt{None} is falsy, but it is not the same object as \texttt{False}.

\begin{verbatim}
print(f"None == False: {None == False}")
print(f"None is False: {None is False}")

print()
print(f"type(None):  {type(None)}")
print(f"type(False): {type(False)}")
\end{verbatim}

Possible output:

\begin{verbatim}
None == False: False
None is False: False

type(None):  <class 'NoneType'>
type(False): <class 'bool'>
\end{verbatim}

\subsubsection{\texttt{None} as a Singleton Object, Not a Null Pointer}

In C, a null pointer is a pointer value that points to no valid object. Python's \texttt{None} should not be understood that way.

\texttt{None} is a real Python object. Names can be bound to it, functions can return it, it has a type, and it has identity.

\begin{verbatim}
a = None
b = None

print(f"a: {a}")
print(f"b: {b}")

print()
print(f"a == b: {a == b}")
print(f"a is b: {a is b}")

print()
print(f"id(a): {id(a)}")
print(f"id(b): {id(b)}")
\end{verbatim}

Possible output:

\begin{verbatim}
a: None
b: None

a == b: True
a is b: True

id(a): 140734431489696
id(b): 140734431489696
\end{verbatim}

There is only one normal \texttt{None} object in a Python process. It is a singleton.

Conceptually:

\begin{verbatim}
a ----\
      ---> None singleton object
b ----/
\end{verbatim}

This is different from an unbound name.

\begin{verbatim}
value = None

print(value)
print(type(value))
\end{verbatim}

This works because the name \texttt{value} exists and points to \texttt{None}.

But this fails:

\begin{verbatim}
try:
    print(missing_value)
except NameError as err:
    print("name lookup failed")
    print(f"message: {err}")
\end{verbatim}

Possible output:

\begin{verbatim}
name lookup failed
message: name 'missing_value' is not defined
\end{verbatim}

The difference is:

\begin{verbatim}
value = None       -> the name exists and refers to None
missing_value      -> the name does not exist here
\end{verbatim}

So \texttt{None} is not ``no variable''. It is an object used to represent no useful value.

\subsubsection{Absence, Default Return Values, and Missing-Value Signaling}

\texttt{None} is commonly used when a value is absent.

\begin{verbatim}
middle_name = None

print(f"middle_name: {middle_name}")
print(f"type:        {type(middle_name)}")

if middle_name is None:
    print("No middle name was provided.")
else:
    print(f"Middle name: {middle_name}")
\end{verbatim}

Possible output:

\begin{verbatim}
middle_name: None
type:        <class 'NoneType'>
No middle name was provided.
\end{verbatim}

This is different from an empty string:

\begin{verbatim}
name1 = None
name2 = ""

print(f"name1: {name1!r}, type={type(name1)}")
print(f"name2: {name2!r}, type={type(name2)}")

print()
print(f"name1 == name2: {name1 == name2}")
print(f"bool(name1): {bool(name1)}")
print(f"bool(name2): {bool(name2)}")
\end{verbatim}

Possible output:

\begin{verbatim}
name1: None, type=<class 'NoneType'>
name2: '', type=<class 'str'>

name1 == name2: False
bool(name1): False
bool(name2): False
\end{verbatim}

Both \texttt{None} and the empty string are falsy, but they do not mean the same thing.

\begin{itemize}
    \item \texttt{None} usually means: no value was supplied, found, or produced.
    \item \texttt{""} means: a string exists, but it contains no characters.
\end{itemize}

A function that does not explicitly return a value returns \texttt{None} automatically.

\begin{verbatim}
def say_quote():
    print("And now for something completely different.")

result = say_quote()

print()
print(f"result: {result}")
print(f"type(result): {type(result)}")
print(f"result is None: {result is None}")
\end{verbatim}

Possible output:

\begin{verbatim}
And now for something completely different.

result: None
type(result): <class 'NoneType'>
result is None: True
\end{verbatim}

The function printed text, but it did not return a useful value. Therefore, Python returned \texttt{None}.

Compare that with a function that returns a value:

\begin{verbatim}
def get_answer():
    return 42

result = get_answer()

print(f"result: {result}")
print(f"type(result): {type(result)}")
print(f"result is None: {result is None}")
\end{verbatim}

Possible output:

\begin{verbatim}
result: 42
type(result): <class 'int'>
result is None: False
\end{verbatim}

\texttt{None} is also useful as a default marker:

\begin{verbatim}
user_name = None

if user_name is None:
    display_name = "Anonymous Pythonista"
else:
    display_name = user_name

print(f"display_name: {display_name}")
\end{verbatim}

Possible output:

\begin{verbatim}
display_name: Anonymous Pythonista
\end{verbatim}

This pattern is common because \texttt{None} clearly separates ``missing value'' from valid values that merely happen to be falsy.

For example, zero may be a valid value:

\begin{verbatim}
score = 0

if score is None:
    print("No score was recorded.")
else:
    print(f"Recorded score: {score}")
\end{verbatim}

Possible output:

\begin{verbatim}
Recorded score: 0
\end{verbatim}

If the program used only a truth test, it would confuse \texttt{0} with absence:

\begin{verbatim}
score = 0

if score:
    print(f"Recorded score: {score}")
else:
    print("This branch also runs for zero.")
\end{verbatim}

Possible output:

\begin{verbatim}
This branch also runs for zero.
\end{verbatim}

So \texttt{None} is the better sentinel when absence must be distinguished from falsy but valid values.

\subsubsection{Identity Testing with \texttt{is None} and \texttt{is not None}}

The standard way to test for \texttt{None} is identity testing:

\begin{verbatim}
value is None
value is not None
\end{verbatim}

Because \texttt{None} is a singleton, identity is the precise question. We are asking whether the name refers to the one \texttt{None} object.

\begin{verbatim}
values = [None, 0, "", False, 42, "D'oh!"]

for value in values:
    print(f"{repr(value):<8} "
          f"is None: {value is None!s:<5} "
          f"is not None: {value is not None!s:<5}")
\end{verbatim}

Possible output:

\begin{verbatim}
None     is None: True  is not None: False
0        is None: False is not None: True 
''       is None: False is not None: True 
False    is None: False is not None: True 
42       is None: False is not None: True 
"D'oh!"  is None: False is not None: True 
\end{verbatim}

This keeps the distinction clean:

\begin{verbatim}
if value is None:
    print("missing")
else:
    print("present")
\end{verbatim}

Do not write this for ordinary \texttt{None} checks:

\begin{verbatim}
if value == None:
    print("missing")
\end{verbatim}

It often works, but it asks the wrong question. The operator \texttt{==} asks whether two objects compare as equal. The operator \texttt{is} asks whether two names refer to the same object.

For \texttt{None}, identity is the intended test.

A compact comparison:

\begin{verbatim}
value = None

print(f"value == None: {value == None}")
print(f"value is None: {value is None}")

value = 0

print()
print(f"value == None: {value == None}")
print(f"value is None: {value is None}")
\end{verbatim}

Possible output:

\begin{verbatim}
value == None: True
value is None: True

value == None: False
value is None: False
\end{verbatim}

The first case looks the same, but the meaning differs. \texttt{is None} is clearer and more exact.

The opposite test is \texttt{is not None}:

\begin{verbatim}
quote = "No soup for you."

if quote is not None:
    print(f"quote exists: {quote}")
else:
    print("quote is missing")
\end{verbatim}

Possible output:

\begin{verbatim}
quote exists: No soup for you.
\end{verbatim}

This does not test whether the string is empty. It tests whether the value is missing.

\begin{verbatim}
quote = ""

if quote is not None:
    print("quote was supplied")
else:
    print("quote is missing")

if quote:
    print("quote has visible content")
else:
    print("quote is empty or otherwise falsy")
\end{verbatim}

Possible output:

\begin{verbatim}
quote was supplied
quote is empty or otherwise falsy
\end{verbatim}

The two tests answer different questions:

\begin{itemize}
    \item \texttt{quote is not None}: was a value supplied?
    \item \texttt{bool(quote)} or \texttt{if quote}: is the supplied value truthy?
\end{itemize}

The practical summary is:

\begin{quote}
\texttt{None} is Python's singleton object for absence or no useful result. It is not a null pointer and not the same thing as \texttt{False}, zero, or an empty string. Test it with \texttt{is None} and \texttt{is not None}.
\end{quote}